% GENERATED_BY: oc_core_monograph_translation_draft_pipeline_v1.py
% AI_ASSISTED_TRANSLATION_DRAFT: TRUE
% THEOREM_REVIEW_STATUS: PENDING
% THEOREM_REVIEW_MODE: FORMAL_AI_GATE__CODEX_CLI_CHATGPT
% THEOREM_REVIEW_REVIEWER_ID: 
% THEOREM_REVIEW_AUTHORITY_CLASS: 
% THEOREM_REVIEWED_AT: 
% THEOREM_REVIEW_DOSSIER_REF: 
% SOURCE_LANGUAGE: EN
% TARGET_LANGUAGE: DE
% SOURCE_EN_REF: content/18_oc_core_1_3_source_audit.tex
\section{Quellenzeugenschaft, Theorem-Schicksal und Publikationsgrenze}
\label{sec:oc13-publication-boundary}

Dieses Kapitel fixiert die Publikationsgrenze von Core~1.3. Seine Aufgabe ist
es, festzulegen, welche wissenschaftliche Autorität kanonisch ist, wie
Quellenzeugenschaft durchgesetzt wird, wie sich das Theorem-Schicksal über die
äußeren Artefakte hinweg trennt und wie spätere Kanäle mit dem Quellkorpus,
SPOT und der Flaggschiff-Monographie ausgerichtet bleiben müssen. Es geht
nicht darum, die formale Lehre oder das Theorem-Gefüge erneut zu begründen;
diese Lasten gehören in die umliegenden Kapitel und Anhänge.

\subsection{Kanonische wissenschaftliche Autorität}

Die kanonische wissenschaftliche Autorität der Veröffentlichung ist die
SPOT-Projektion
\occode{releases/oc_core_1_3/editorial/OC_CORE_1_3_SCIENCE_SPOT_latest.json}
(surface id \occode{OC_CORE_1_3_SCIENCE_SPOT}; current SHA recorded in its
\occode{metadata.repo_sha} field). Das quellenbesitzende wissenschaftliche
Korpus unter
\occode{releases/oc_core_1_3/editorial/science_sources/} ist der Ort des
Verfassens und der Zeugenschaft, aus dem diese SPOT-Projektion abgeleitet wird.
Historische Quellenzeugnisse, einschließlich früherer Core~1.2/Core~1.3-
Reparatur-Baselines, bleiben Provenienzanker; sie sind nicht die aktuelle
Projektionsautorität. Die gegenwärtige Master-Monographie ist die
reader-facing Flaggschiff-Projektion und der zusammengestellte Referenzband,
der aus diesen Aufzeichnungen abgeleitet ist. Sie hält Orientierungsskala,
formale Lehre, Theorem-Pfad, empirische Disziplin, Synthese, praktische
Nutzbarkeit und Audit-Lagen gemeinsam sichtbar, verdrängt aber SPOT nicht als
endgültige Autorität für Anspruchsumfang, Promotionsstatus oder Beweisgrenze.

Diese Regel ist wichtig, weil grundlagentheoretische Programme im Zeitverlauf
mehrere unterschiedliche Textschichten ansammeln: quellennahe Definitionen und
Theoreme, erläuternde Prosa, Reparaturvermerke, kanalspezifische
Zusammenfassungen und spätere operative Pakete. Core~1.3 behandelt diese
Schichten als unterscheidbar, nicht als austauschbar.

\subsection{Quellenzeugenschaft}

Core~1.3 ist streng quellenfirst im Sinn, dass die im öffentlichen Release
verwendeten Flaggschiff-Ansprüche an öffentliche Quellzeugnisse gebunden
bleiben müssen und nicht allein an späte Zusammenfassungsprosa. Der Punkt ist
nicht archivalischer Fetischismus. Der Punkt ist, die Drift zwischen fünf
Dingen zu verhindern, die in großen Forschungsprogrammen oft ineinander
verwaschen:

\begin{itemize}
\item theoremtragende Aussagen,
\item strukturelle und interpretative Erläuterung,
\item Reparatur- und Release-Notizen,
\item verdichtete öffentliche Pakete,
\item und spätere operative Projektionen.
\end{itemize}

Wenn diese Grenze verschwimmt, kann ein Zusammenfassungssatz beginnen, sich als
das Theorem selbst auszugeben, oder ein Release-Dashboard kann das eigentliche
wissenschaftliche Protokoll zu ersetzen beginnen. Core~1.3 verhindert diesen
Fehler, indem es die Quellenzeugenschaft an das quellenbesitzende
wissenschaftliche Korpus und die kanonische SPOT-Projektion bindet, statt sie
in verborgene Steuerdateien auszulagern.

\subsection{Theorem-Schicksal und Umfangsdisziplin}

Auch die vorliegende Veröffentlichung trennt das Theorem-Schicksal explizit.
Nicht jede geerbte Zeile erfüllt in jedem äußeren Artefakt dieselbe Rolle.

\begin{itemize}
\item \occode{THEOREM_NATIVE}-Zeilen dürfen positive theorem-native
      Aussagen tragen, wenn ihre Beweis- und Abhängigkeitsroute exakt bleibt.
\item \occode{BRIDGE_ONLY}-Zeilen können unter expliziten Umfangsnotizen
      wissenschaftlich nützlich bleiben, dürfen aber nicht als theorem-native
      aufgestuft werden.
\item \occode{FRAME_ONLY}- und \occode{EXTENSION}-Zeilen gehören in
      Stützschichten, Anhänge, engere Begleitartefakte oder künftige Arbeit
      statt in den positiven Kern jedes äußeren Papiers.
\item \occode{REFUTED}-Zeilen sind vom positiven Beweisweg ausgeschlossen.
      Sie dürfen nur als negative Audit-Historie oder als Falsifikator-Record
      erwähnt werden, nicht als lebende Unterstützung.
\item Frontier- oder Hypothesenmaterial darf nur unter sichtbar markierter
      Frontier-Sprache diskutiert werden und darf nicht durch
      Zusammenfassungsprosa stillschweigend aufgestuft werden.
\end{itemize}

Der öffentliche positive Schlussbehauptungsraum ist genau auf
\occode{CORE_1_3_SCIENCE_ONLY} beschränkt. \occode{FRAME_ONLY}-Zeilen können
innerhalb der Syntheseoberfläche weiterhin als Frame-Level-\texttt{PASS}-Zeilen
schließen, werden dadurch aber nicht zu eigenständigen theorem-native
empirischen Promotions. \occode{EXTENSION}- und \occode{REFUTED}-Zeilen
bleiben aus unterschiedlichen Gründen außerhalb des öffentlichen
\texttt{PASS}-Umfangs: \occode{EXTENSION}-Zeilen sind Front-/Support-Lanes, die
auf spätere Promotion oder Ablehnung warten, während
\occode{REFUTED}-Zeilen negative Audit-Lanes sind. Keine der beiden Klassen
kann in dieser Veröffentlichung positive Schlussbehauptungen tragen.
Zur Vermeidung von Missverständnissen ist eine \occode{FRAME_ONLY}
\occode{PASS}-Zeile vom öffentlichen positiven Schlussbehauptungsraum
ausgeschlossen, sofern sie nicht in einer späteren Veröffentlichung mit der
erforderlichen theorem-native Unterstützung in \occode{CORE_1_3_SCIENCE_ONLY}
aufgestuft wird.

Das ist kein Rückzug. Es ist eine Bedingung ehrlicher Veröffentlichung. Die
Monographie kann groß bleiben, ohne rhetorisch aufzublähen, nur wenn die
Leserschaft erkennen kann, welche Zeilen theoremtragend, welche begrenzt und
welche Stütz- oder Frontier-Material sind.

\subsection{Bezug zu Journal-Kern und Begleitartefakten}

Core~1.3 übernimmt daher eine einseitige, reader-facing Publikationslogik:
\[
\begin{gathered}
\text{quellenbesitzendes wissenschaftliches Korpus}\\
\text{plus kanonisches SPOT}\\
\Downarrow\\
\text{Master-Monographie}\\
\Downarrow\\
\text{abgeleitete äußere Artefakte}
\end{gathered}
\]
Dieses Diagramm beschreibt die Publikationsreihenfolge, nicht den gesamten
Provenienzgraphen des Release-Pakets. Mehrere redaktionelle JSON-Oberflächen,
generierte Wissenschaftskapitel und Release-Matrizen sind direkte
SPOT-Projektionen oder Spiegel anderer maßgeblicher kanonischer Oberflächen.
Sie werden nicht dadurch wahr, dass die Master-Monographie sie zitiert, und
sie sind nicht allein Nachfahren der Monographie. Die Monographie ist das
Flaggschiff-Leseartefakt; das quellenbesitzende wissenschaftliche Korpus bleibt
der Ort des Verfassens und der Zeugenschaft, und SPOT bleibt die kanonische
Projektionsautorität für Release-Oberflächen.

Der Journal-Core-Artikel ist bewusst enger gefasst. Begleit-Dossiers,
Matrizen und Release-Oberflächen sind bewusst enger oder operativer gefasst.
Diese Artefakte dürfen das Flaggschiff-Material für einen bestimmten Zweck
komprimieren, auswählen oder umorganisieren, aber sie dürfen den
wissenschaftlichen Anspruch nicht vergrößern, die Notation nicht ausweiten und
die hier festgelegte Theorem-Schicksalsdisziplin nicht überholen.

Die praktische Konsequenz ist einfach. Wenn ein äußeres,
reader-facing Artefakt etwas Stärkeres behauptet, als Quellkorpus und SPOT
autorisiert haben, dann ist das äußere Artefakt falsch, auch wenn es kürzer,
sauberer oder leichter vermarktbar ist. Die Flaggschiff-Grenze ist die
Publikationsregel, die diese Obergrenze durchsetzt; sie ist keine zusätzliche
wissenschaftliche Autorität.

\subsection{Downstream-Projektionsregel}

Dieselbe Grenze stabilisiert auch die spätere Verwendung in Review,
Zusammenarbeit, Produkt, Business und öffentlicher Kommunikation.
Verschiedene Kanäle dürfen legitimerweise verschiedene Teile des
Flaggschiff-Bandes zitieren:

\begin{itemize}
\item den Theorem-Pfad, wenn begrenzter Beweis das Thema ist,
\item den empirischen Pfad, wenn Promotionsdisziplin das Thema ist,
\item den praktischen Pfad, wenn Nutzung und Vergleich das Thema sind,
\item die Audit-Anhänge, wenn Anfechtung und Nachverfolgbarkeit das Thema sind.
\end{itemize}

Was sie nicht dürfen, ist eine zweite wissenschaftliche Autorität zu erfinden.
Quellenzeugenschaft, Theorem-Schicksal und Publikationsgrenze existieren
deshalb, damit die externe Nutzung knapp bleiben kann, ohne interpretativ
gesetzlos zu werden.

\paragraph{Brücke zum nächsten Kapitel.}
Mit der fixierten Publikationsgrenze ist die nächste Aufgabe die innere
Konsistenz: welche Aussagenklassen der Flaggschiff-Band verwendet, warum die
Hierarchie durch \(K_{12}\) läuft und wie die oberen Ebenen gesetzmäßig
bleiben, ohne rhetorisch aufzublähen.

\clearpage

\subsection{External criticism closure route}
\label{sec:oc13-source-audit-external-criticism-closure}

The Stanislav Tsukrov external review package is closed through Appendix~\ref{sec:oc13-external-criticism-closure}. The route records which objections received evidence boundaries, which became proof tasks, and which duplicate revisions were merged. It does not move private review text into the public release and it does not promote a canonical claim.

% SOURCE_EN_REF: content/18_oc_core_1_3_source_audit.tex

\subsection{Source-tradition remediation after external review}
\label{sec:oc13-source-tradition-remediation-v3}

The Stanislav Tsukrov review identified a real scholarly weakness: Core~1.3
already cited physics-adjacent complexity literature, but it did not adequately
name the systems, autopoiesis, identity, complexity-measure, tipping-point, and
formal-geometry traditions in which several OC constructs must be judged. This
release therefore lowers the novelty ceiling and makes the inherited ground
explicit.

The hierarchy and systems-language route is now read against general systems
theory and mathematical systems theory, including von Bertalanffy,
Boulding, Miller, Mesarovic--Takahara, and Klir
\cite{bertalanffy_general_system_theory,boulding_general_systems_theory,miller_living_systems,mesarovic_takahara_general_systems,klir_architecture_systems}.
The rebirth and post-collapse identity route is read against autopoiesis and
identity/persistence debates, including Maturana--Varela, Varela, Parfit, and
Wimsatt \cite{maturana_varela_autopoiesis,varela_biological_autonomy,parfit_reasons_persons,wimsatt_reductionism_levels}.
The complexity and early-warning route is read against effective complexity,
predictive information, integrated information, and critical-transition
literature \cite{gell_mann_effective_complexity,bialek_predictability_complexity,tononi_information_integration,scheffer_critical_transitions,dakos_early_warning}.
Formal differential-geometric machinery is treated as inherited mathematical
tooling rather than as an OC novelty claim; Olver is the current compact source
witness for the jet/Lie-group side of that boundary \cite{olver_lie_groups}.

The repaired claim is narrow: OC contributes a typed cross-level interface,
collapse/residue/rebirth grammar, and release discipline. It does not claim to
replace GST, autopoiesis, complexity theory, identity metaphysics, tipping-point
analysis, or formal differential geometry.
